% Options for packages loaded elsewhere
\PassOptionsToPackage{unicode}{hyperref}
\PassOptionsToPackage{hyphens}{url}
%
\documentclass[
]{article}
\usepackage{amsmath,amssymb}
\usepackage{iftex}
\ifPDFTeX
  \usepackage[T1]{fontenc}
  \usepackage[utf8]{inputenc}
  \usepackage{textcomp} % provide euro and other symbols
\else % if luatex or xetex
  \usepackage{unicode-math} % this also loads fontspec
  \defaultfontfeatures{Scale=MatchLowercase}
  \defaultfontfeatures[\rmfamily]{Ligatures=TeX,Scale=1}
\fi
\usepackage{lmodern}
\ifPDFTeX\else
  % xetex/luatex font selection
\fi
% Use upquote if available, for straight quotes in verbatim environments
\IfFileExists{upquote.sty}{\usepackage{upquote}}{}
\IfFileExists{microtype.sty}{% use microtype if available
  \usepackage[]{microtype}
  \UseMicrotypeSet[protrusion]{basicmath} % disable protrusion for tt fonts
}{}
\makeatletter
\@ifundefined{KOMAClassName}{% if non-KOMA class
  \IfFileExists{parskip.sty}{%
    \usepackage{parskip}
  }{% else
    \setlength{\parindent}{0pt}
    \setlength{\parskip}{6pt plus 2pt minus 1pt}}
}{% if KOMA class
  \KOMAoptions{parskip=half}}
\makeatother
\usepackage{xcolor}
\usepackage{graphicx}
\makeatletter
\def\maxwidth{\ifdim\Gin@nat@width>\linewidth\linewidth\else\Gin@nat@width\fi}
\def\maxheight{\ifdim\Gin@nat@height>\textheight\textheight\else\Gin@nat@height\fi}
\makeatother
% Scale images if necessary, so that they will not overflow the page
% margins by default, and it is still possible to overwrite the defaults
% using explicit options in \includegraphics[width, height, ...]{}
\setkeys{Gin}{width=\maxwidth,height=\maxheight,keepaspectratio}
% Set default figure placement to htbp
\makeatletter
\def\fps@figure{htbp}
\makeatother
\setlength{\emergencystretch}{3em} % prevent overfull lines
\providecommand{\tightlist}{%
  \setlength{\itemsep}{0pt}\setlength{\parskip}{0pt}}
\setcounter{secnumdepth}{-\maxdimen} % remove section numbering
\ifLuaTeX
  \usepackage{selnolig}  % disable illegal ligatures
\fi
\IfFileExists{bookmark.sty}{\usepackage{bookmark}}{\usepackage{hyperref}}
\IfFileExists{xurl.sty}{\usepackage{xurl}}{} % add URL line breaks if available
\urlstyle{same}
\hypersetup{
  pdftitle={Coursework 1},
  hidelinks,
  pdfcreator={LaTeX via pandoc}}

\title{Coursework 1}
\author{}
\date{}

\begin{document}
\maketitle

a)
\includegraphics{C:/Users/artem/OneDrive/Documents/Obsidian Vault/Pasted image 20231031202306.png}\\
b)

\begin{verbatim}
pothole list = [n1... nk]                       O(1)
distance covered = 0                            O(1)
i = 1                                           O(1)
While i < size of pothole list:                 O(n)
    current pothole = pothole list[i]           O(1)
    if current pothole > distance covered:
        Place a slab with the left end just covering the pothole                                         O(1)
        distance covered = current pothole + 1  O(1)
    i = i + 1                                   O(1)
Copy
\end{verbatim}

\textbf{Proof}:\\
Given that every slab has a pothole at its left edge, a suboptimal
covering of potholes happens in the following cases:\\
1. A slab is placed that doesn\textquotesingle t cover any potholes\\
2. 2 slabs cover a range of potholes that are at most 1m apart (They
could be covered by just 1 slab)\\
My algorithm doesn\textquotesingle t place a slab in either of these
cases:\\
1. My algorithm places slabs only when a pothole is present, placing the
left edge of the slab on the pothole\\
2. Slab n only placed if slab n-1 is placed is at least 1m away, and
since slab n-1 must have a pothole at the left edge of it, the potholes
covered by slab n must be at least 1m away from this very leftmost
pothole under slab n-1

Since this basically just iterates through the list of potholes which
takes O(n) time, with all other operations being O(1), the algorithm
runs in O(n) time

c)\\
\includegraphics{C:/Users/artem/OneDrive/Documents/Obsidian Vault/Pasted image 20231031202808.png}\\
6 Slabs

d)

\begin{enumerate}
\tightlist
\item
  It doesn\textquotesingle t ensure the condition for slab{[}i{]} to be
  the same as slab{[}i+2{]}
\item
  Doing this greedily, like how the AI has done it,
  doesn\textquotesingle t always return an optimal number of
  operations.\\
  For example, if the 3...m slabs were already in an optimal order but
  the first 2 weren\textquotesingle t, this algorithm would change the
  color of the first 2 and potentially cause a chain reaction causing it
  to change the color of all the rest of the slabs, unnecessarily
\end{enumerate}

e)\\
Starting arrangement: {[}R, G, B, G, R, Y, P{]}\\
Final arrangement: {[}R, G, R, G, R, G, R{]}

Count the number of each coloured tile in the even indexes and the odd
ones. \emph{O(n)}\\
Choose pair of tile colours, one from each list, whose tallies add up to
the highest number, but which aren\textquotesingle t the same colour (so
choose second best pair if both same colour). \emph{O(nlogn)}

Optimal:

Final arrangement of the algorithm colours the tiles is such that:

\begin{itemize}
\tightlist
\item
  There is only one colour in the even tiles

  \begin{itemize}
  \tightlist
  \item
    due to rule that slab{[}i{]} = slab{[}i+2{]}
  \end{itemize}
\item
  There is a different colour in the odd tiles

  \begin{itemize}
  \tightlist
  \item
    due to rule that slab{[}i{]} =/= slab{[}i-1{]}
  \end{itemize}
\end{itemize}

An optimal algorithm colours the slabs in this way, and for the even
colour repaints the least amount of slabs to this colour, and same for
odd slabs/colour.\\
This means that the optimal colour we choose to repaint the even slabs
should be the colour that appears most often in the even spaces in the
original arrangement, and same for the odd slabs.

Since my algorithm counts the number of slabs painted in each colour in
the even and odd spaces, and maximises the combination of the two, while
keeping completeness by only choosing them only if they are different
colours, it must be optimal.

Time complexity = O(n) + O(nlogn) = O(nlogn)

f)\\
The algorithm takes an arbitrary array as an input and outputs an
ordered array with specific conditions to the ordering.\\
For any algorithm that transforms one such array of elements of size n
into another such array, the fastest it can possibly be is O(n) as it at
the least it needs to iterate through the array in order to check that
it follows the desired conditions.

\end{document}
